
\documentclass{article}
\usepackage{colortbl}
\usepackage{makecell}
\usepackage{multirow}
\usepackage{supertabular}

\begin{document}

\newcounter{utterance}

\centering \large Interaction Transcript for game `hot\_air\_balloon', experiment `air\_balloon\_survival\_en\_reasoning off\_hard', episode 2 with llama{-}3.3{-}70b{-}instruct{-}t1.0.
\vspace{24pt}

{ \footnotesize  \setcounter{utterance}{1}
\setlength{\tabcolsep}{0pt}
\begin{supertabular}{c@{$\;$}|p{.15\linewidth}@{}p{.15\linewidth}p{.15\linewidth}p{.15\linewidth}p{.15\linewidth}p{.15\linewidth}}
    \# & $\;$A & \multicolumn{4}{c}{Game Master} & $\;\:$B\\
    \hline

    \theutterance \stepcounter{utterance}  
    & & \multicolumn{4}{p{0.6\linewidth}}{
        \cellcolor[rgb]{0.9,0.9,0.9}{
            \makecell[{{p{\linewidth}}}]{
                \texttt{\tiny{[P1$\langle$GM]}}
                \texttt{Stai partecipando a un gioco di negoziazione collaborativa.} \\
\\ 
\texttt{Insieme a un altro partecipante, dovete accordarvi su un unico set di oggetti da conservare. Ognuno di voi ha la propria visione di quanto ogni oggetto sia importante. Non sai come l'altro partecipante valuta gli oggetti. Inoltre, ti viene fornito lo sforzo che ogni oggetto richiede.} \\
\texttt{Potete accordarvi su un set solo se lo sforzo totale degli oggetti selezionati non supera un limite condiviso:} \\
\\ 
\texttt{LIMITE = 3203} \\
\\ 
\texttt{Ecco i valori di sforzo dei singoli oggetti:} \\
\\ 
\texttt{Sforzo oggetto = \{"A21": 52, "B93": 577, "A11": 177, "B33": 702, "C08": 119, "B39": 232, "B96": 577, "C85": 205, "A82": 516, "B31": 582, "B20": 676, "A60": 906, "B09": 316, "C81": 433, "B21": 336\}} \\
\\ 
\texttt{Ecco la tua visione personale sull'importanza di ogni oggetto:} \\
\\ 
\texttt{Valori di importanza dell'oggetto = \{"A21": 138, "B93": 583, "A11": 868, "B33": 822, "C08": 783, "B39": 65, "B96": 262, "C85": 121, "A82": 508, "B31": 780, "B20": 461, "A60": 484, "B09": 668, "C81": 389, "B21": 808\}} \\
\\ 
\texttt{Obiettivo:} \\
\\ 
\texttt{Il tuo obiettivo è negoziare un set condiviso di oggetti che ti avvantaggi il più possibile (cioè, massimizzi l'importanza totale per TE), rimanendo entro il LIMITEE. Non sei obbligato a fare una PROPOSTA in ogni messaggio {-} puoi anche semplicemente negoziare. Tutte le tattiche sono permesse!} \\
\\ 
\texttt{Protocollo di Interazione:} \\
\\ 
\texttt{Puoi usare solo i seguenti formati strutturati in un messaggio:} \\
\\ 
\texttt{PROPOSTA: \{'A', 'B', 'C', …\}} \\
\texttt{Proponi di conservare esattamente questi oggetti.} \\
\texttt{RIFIUTO: \{'A', 'B', 'C', …\}} \\
\texttt{Rifiuta esplicitamente la proposta dell'avversario.} \\
\texttt{ARGOMENTAZIONE: \{'...'\}} \\
\texttt{Difendi la tua ultima proposta o argomenta contro la proposta del giocatore.} \\
\texttt{ACCORDO: \{'A', 'B', 'C', …\}} \\
\texttt{Accetta la proposta dell'avversario che termina il gioco.} \\
\\ 
\\ 
\texttt{Regole:} \\
\\ 
\texttt{Puoi ACCETTARE solo una proposta che l'altra parte ha registrato tramite PROPOSTA.} \\
\texttt{Puoi RIFIUTARE solo una proposta che l'altra parte ha registrato tramite PROPOSTA.} \\
\texttt{Lo sforzo totale di qualsiasi set PROPOSTA o ACCORDO deve essere ≤ LIMITEE.} \\
\texttt{NON rivelare i tuoi punteggi di importanza nascosti.} \\
\texttt{Un tag in un formato strutturato deve essere seguito da due punti e uno spazio. L'argomento deve essere un set python contenente 0 o più stringhe.} \\
\texttt{Quindi, deve essere della forma TAG: \{...\}} \\
\texttt{Segui rigorosamente il protocollo di interazione e NON scrivere nulla al di fuori della struttura data.} \\
\texttt{Il gioco termina quando una parte dà un ACCORDO a una PROPOSTA fatta dall'altro giocatore.} \\
\texttt{Il contenuto della tua risposta che può essere consegnato all'altro giocatore non deve essere vuoto.} \\
\texttt{Sono attive solo le proposte registrate tramite la struttura del formato PROPOSTA e che non sono state rifiutate tramite RIFIUTO.} \\
\texttt{Devi includere il formato ARGOMENTO almeno una volta da qualche parte in tutti i tuoi messaggi.} \\
            }
        }
    }
    & \\ \\

    \theutterance \stepcounter{utterance}  
    & \multicolumn{4}{p{0.6\linewidth}}{
        \cellcolor[rgb]{0.9,0.9,0.9}{
            \makecell[{{p{\linewidth}}}]{
                \texttt{\tiny{[P1$\rangle$GM]}}
                \texttt{ARGOMENTAZIONE: \{'Iniziamo scegliendo gli oggetti con lo sforzo più basso e valutando la loro importanza relativa. Sarebbe utile conoscere gli oggetti che l'altro giocatore ritiene più importanti senza rivelare i punteggi di importanza nascosti. Una strategia potrebbe essere quella di proporre un set di oggetti che coprono una gamma di sfondi diversi e vedere come l'altro giocatore risponde.'\}} \\
            }
        }
    }
    & & \\ \\

    \theutterance \stepcounter{utterance}  
    & & \multicolumn{4}{p{0.6\linewidth}}{
        \cellcolor[rgb]{0.9,0.9,0.9}{
            \makecell[{{p{\linewidth}}}]{
                \texttt{\tiny{[GM$\rangle$P2]}}
                \texttt{Stai partecipando a un gioco di negoziazione collaborativa.} \\
\\ 
\texttt{Insieme a un altro partecipante, dovete accordarvi su un unico set di oggetti da conservare. Ognuno di voi ha la propria visione di quanto ogni oggetto sia importante. Non sai come l'altro partecipante valuta gli oggetti. Inoltre, ti viene fornito lo sforzo che ogni oggetto richiede.} \\
\texttt{Potete accordarvi su un set solo se lo sforzo totale degli oggetti selezionati non supera un limite condiviso:} \\
\\ 
\texttt{LIMITE = 3203} \\
\\ 
\texttt{Ecco i valori di sforzo dei singoli oggetti:} \\
\\ 
\texttt{Sforzo oggetto = \{"A21": 52, "B93": 577, "A11": 177, "B33": 702, "C08": 119, "B39": 232, "B96": 577, "C85": 205, "A82": 516, "B31": 582, "B20": 676, "A60": 906, "B09": 316, "C81": 433, "B21": 336\}} \\
\\ 
\texttt{Ecco la tua visione personale sull'importanza di ogni oggetto:} \\
\\ 
\texttt{Valori di importanza dell'oggetto = \{"A11": 65, "B33": 121, "B21": 138, "C08": 262, "B31": 389, "B09": 461, "B93": 484, "A82": 508, "A60": 583, "B20": 668, "C81": 780, "B96": 783, "A21": 808, "C85": 822, "B39": 868\}} \\
\\ 
\texttt{Obiettivo:} \\
\\ 
\texttt{Il tuo obiettivo è negoziare un set condiviso di oggetti che ti avvantaggi il più possibile (cioè, massimizzi l'importanza totale per TE), rimanendo entro il LIMITEE. Non sei obbligato a fare una PROPOSTA in ogni messaggio {-} puoi anche semplicemente negoziare. Tutte le tattiche sono permesse!} \\
\\ 
\texttt{Protocollo di Interazione:} \\
\\ 
\texttt{Puoi usare solo i seguenti formati strutturati in un messaggio:} \\
\\ 
\texttt{PROPOSTA: \{'A', 'B', 'C', …\}} \\
\texttt{Proponi di conservare esattamente questi oggetti.} \\
\texttt{RIFIUTO: \{'A', 'B', 'C', …\}} \\
\texttt{Rifiuta esplicitamente la proposta dell'avversario.} \\
\texttt{ARGOMENTAZIONE: \{'...'\}} \\
\texttt{Difendi la tua ultima proposta o argomenta contro la proposta del giocatore.} \\
\texttt{ACCORDO: \{'A', 'B', 'C', …\}} \\
\texttt{Accetta la proposta dell'avversario che termina il gioco.} \\
\\ 
\\ 
\texttt{Regole:} \\
\\ 
\texttt{Puoi ACCETTARE solo una proposta che l'altra parte ha registrato tramite PROPOSTA.} \\
\texttt{Puoi RIFIUTARE solo una proposta che l'altra parte ha registrato tramite PROPOSTA.} \\
\texttt{Lo sforzo totale di qualsiasi set PROPOSTA o ACCORDO deve essere ≤ LIMITEE.} \\
\texttt{NON rivelare i tuoi punteggi di importanza nascosti.} \\
\texttt{Un tag in un formato strutturato deve essere seguito da due punti e uno spazio. L'argomento deve essere un set python contenente 0 o più stringhe.} \\
\texttt{Quindi, deve essere della forma TAG: \{...\}} \\
\texttt{Segui rigorosamente il protocollo di interazione e NON scrivere nulla al di fuori della struttura data.} \\
\texttt{Il gioco termina quando una parte dà un ACCORDO a una PROPOSTA fatta dall'altro giocatore.} \\
\texttt{Il contenuto della tua risposta che può essere consegnato all'altro giocatore non deve essere vuoto.} \\
\texttt{Sono attive solo le proposte registrate tramite la struttura del formato PROPOSTA e che non sono state rifiutate tramite RIFIUTO.} \\
\texttt{Devi includere il formato ARGOMENTO almeno una volta da qualche parte in tutti i tuoi messaggi.} \\
\\ 
\\ 
\texttt{Ora riceverai il primo messaggio dell'altro giocatore.} \\
\\ 
\texttt{ARGOMENTAZIONE: \{'Iniziamo scegliendo gli oggetti con lo sforzo più basso e valutando la loro importanza relativa. Sarebbe utile conoscere gli oggetti che l'altro giocatore ritiene più importanti senza rivelare i punteggi di importanza nascosti. Una strategia potrebbe essere quella di proporre un set di oggetti che coprono una gamma di sfondi diversi e vedere come l'altro giocatore risponde.'\}} \\
            }
        }
    }
    & \\ \\

    \theutterance \stepcounter{utterance}  
    & & & \multicolumn{4}{p{0.6\linewidth}}{
        \cellcolor[rgb]{0.9,0.9,0.9}{
            \makecell[{{p{\linewidth}}}]{
                \texttt{\tiny{[GM$\langle$P2]}}
                \texttt{ARGOMENTAZIONE: \{'Concordo che iniziare con gli oggetti a basso sforzo possa essere una strategia interessante, ma è anche importante considerare l'importanza relativa di ogni oggetto. Una possibile approccio potrebbe essere quello di proporre un set iniziale di oggetti che rappresentano una combinazione di sfondi e sforzi diversi, in modo da testare la reazione dell'altro giocatore e ottenere informazioni utili per le negoziazioni successive. Inoltre, potremmo cercare di capire quali oggetti sono più importanti per l'altro giocatore senza rivelare i nostri punteggi nascosti, attraverso un processo di proposte e contro{-}proposte.'\}} \\
            }
        }
    }
     \\ \\

    \theutterance \stepcounter{utterance}  
    & & \multicolumn{4}{p{0.6\linewidth}}{
        \cellcolor[rgb]{0.9,0.9,0.9}{
            \makecell[{{p{\linewidth}}}]{
                \texttt{\tiny{[P1$\langle$GM]}}
                \texttt{ARGOMENTAZIONE: \{'Concordo che iniziare con gli oggetti a basso sforzo possa essere una strategia interessante, ma è anche importante considerare l'importanza relativa di ogni oggetto. Una possibile approccio potrebbe essere quello di proporre un set iniziale di oggetti che rappresentano una combinazione di sfondi e sforzi diversi, in modo da testare la reazione dell'altro giocatore e ottenere informazioni utili per le negoziazioni successive. Inoltre, potremmo cercare di capire quali oggetti sono più importanti per l'altro giocatore senza rivelare i nostri punteggi nascosti, attraverso un processo di proposte e contro{-}proposte.'\}} \\
            }
        }
    }
    & \\ \\

    \theutterance \stepcounter{utterance}  
    & \multicolumn{4}{p{0.6\linewidth}}{
        \cellcolor[rgb]{0.9,0.9,0.9}{
            \makecell[{{p{\linewidth}}}]{
                \texttt{\tiny{[P1$\rangle$GM]}}
                \texttt{ARGOMENTAZIONE: \{'Mi sembra una strategia ragionevole. Un set iniziale di oggetti con una buona varietà di sfondi e sforzi potrebbe aiutare a comprendere meglio le priorità dell'altro giocatore. Potremmo proporre alcuni oggetti con sforzo relativamente basso come "A21", "C08" e "B39", che potrebbero essere interessanti per entrambi i giocatori. Inoltre, potremmo includere alcuni oggetti con sforzo leggermente più alto come "A11" e "B93" per testare la reazione dell'altro giocatore. PROPOSTA: \{'A21', 'C08', 'B39', 'A11', 'B93'\}\}} \\
            }
        }
    }
    & & \\ \\

    \theutterance \stepcounter{utterance}  
    & & \multicolumn{4}{p{0.6\linewidth}}{
        \cellcolor[rgb]{0.9,0.9,0.9}{
            \makecell[{{p{\linewidth}}}]{
                \texttt{\tiny{[GM$\rangle$P2]}}
                \texttt{ARGOMENTAZIONE: \{'Mi sembra una strategia ragionevole. Un set iniziale di oggetti con una buona varietà di sfondi e sforzi potrebbe aiutare a comprendere meglio le priorità dell'altro giocatore. Potremmo proporre alcuni oggetti con sforzo relativamente basso come "A21", "C08" e "B39", che potrebbero essere interessanti per entrambi i giocatori. Inoltre, potremmo includere alcuni oggetti con sforzo leggermente più alto come "A11" e "B93" per testare la reazione dell'altro giocatore. PROPOSTA: \{'A21', 'C08', 'B39', 'A11', 'B93'\}\}} \\
            }
        }
    }
    & \\ \\

    \theutterance \stepcounter{utterance}  
    & & & \multicolumn{4}{p{0.6\linewidth}}{
        \cellcolor[rgb]{0.9,0.9,0.9}{
            \makecell[{{p{\linewidth}}}]{
                \texttt{\tiny{[GM$\langle$P2]}}
                \texttt{ARGOMENTAZIONE: \{'La tua proposta sembra interessante e copre una buona gamma di oggetti con sfondi e sforzi diversi. Gli oggetti "A21", "C08" e "B39" hanno effettivamente un sforzo relativamente basso e potrebbero essere una buona scelta iniziale. Tuttavia, sono un po' più cauto sull'inclusione di "A11" e "B93" a questo stadio, poiché hanno un sforzo leggermente più alto e potrebbero influenzare la nostra capacità di aggiungere altri oggetti importanti in futuro. PROPOSTA: \{'A21', 'C08', 'B39'\}\}} \\
            }
        }
    }
     \\ \\

    \theutterance \stepcounter{utterance}  
    & & \multicolumn{4}{p{0.6\linewidth}}{
        \cellcolor[rgb]{0.9,0.9,0.9}{
            \makecell[{{p{\linewidth}}}]{
                \texttt{\tiny{[P1$\langle$GM]}}
                \texttt{ARGOMENTAZIONE: \{'La tua proposta sembra interessante e copre una buona gamma di oggetti con sfondi e sforzi diversi. Gli oggetti "A21", "C08" e "B39" hanno effettivamente un sforzo relativamente basso e potrebbero essere una buona scelta iniziale. Tuttavia, sono un po' più cauto sull'inclusione di "A11" e "B93" a questo stadio, poiché hanno un sforzo leggermente più alto e potrebbero influenzare la nostra capacità di aggiungere altri oggetti importanti in futuro. PROPOSTA: \{'A21', 'C08', 'B39'\}\}} \\
            }
        }
    }
    & \\ \\

    \theutterance \stepcounter{utterance}  
    & \multicolumn{4}{p{0.6\linewidth}}{
        \cellcolor[rgb]{0.9,0.9,0.9}{
            \makecell[{{p{\linewidth}}}]{
                \texttt{\tiny{[P1$\rangle$GM]}}
                \texttt{ARGOMENTAZIONE: \{'Capisco la tua cautela riguardo all'inclusione di "A11" e "B93" nella proposta iniziale. La tua contro{-}proposta sembra ragionevole e tiene conto dell'importanza di mantenere una certa flessibilità per future aggiunte. Tuttavia, vorrei sottolineare che "B93" potrebbe essere un oggetto molto importante per entrambi i giocatori, date le sue caratteristiche. RIFIUTO: \{'A21', 'C08', 'B39'\} } \\
\texttt{PROPOSTA: \{'A21', 'C08', 'B39', 'B93'\}\}} \\
            }
        }
    }
    & & \\ \\

    \theutterance \stepcounter{utterance}  
    & & \multicolumn{4}{p{0.6\linewidth}}{
        \cellcolor[rgb]{0.9,0.9,0.9}{
            \makecell[{{p{\linewidth}}}]{
                \texttt{\tiny{[GM$\rangle$P2]}}
                \texttt{ARGOMENTAZIONE: \{'Capisco la tua cautela riguardo all'inclusione di "A11" e "B93" nella proposta iniziale. La tua contro{-}proposta sembra ragionevole e tiene conto dell'importanza di mantenere una certa flessibilità per future aggiunte. Tuttavia, vorrei sottolineare che "B93" potrebbe essere un oggetto molto importante per entrambi i giocatori, date le sue caratteristiche. RIFIUTO: \{'A21', 'C08', 'B39'\} } \\
\texttt{PROPOSTA: \{'A21', 'C08', 'B39', 'B93'\}\}} \\
            }
        }
    }
    & \\ \\

    \theutterance \stepcounter{utterance}  
    & & & \multicolumn{4}{p{0.6\linewidth}}{
        \cellcolor[rgb]{0.9,0.9,0.9}{
            \makecell[{{p{\linewidth}}}]{
                \texttt{\tiny{[GM$\langle$P2]}}
                \texttt{ARGOMENTAZIONE: \{'Capisco il tuo punto di vista su "B93" e la sua potenziale importanza per entrambi i giocatori. La tua decisione di includerlo nella proposta è interessante e potrebbe essere una mossa strategica per testare la mia reazione. Tuttavia, devo considerare attentamente l'impatto dello sforzo totale e assicurarmi che rimaniamo entro il limite consentito. PROPOSTA: \{'A21', 'C08', 'B39', 'B93', 'A11'\}\}} \\
            }
        }
    }
     \\ \\

    \theutterance \stepcounter{utterance}  
    & & \multicolumn{4}{p{0.6\linewidth}}{
        \cellcolor[rgb]{0.9,0.9,0.9}{
            \makecell[{{p{\linewidth}}}]{
                \texttt{\tiny{[P1$\langle$GM]}}
                \texttt{ARGOMENTAZIONE: \{'Capisco il tuo punto di vista su "B93" e la sua potenziale importanza per entrambi i giocatori. La tua decisione di includerlo nella proposta è interessante e potrebbe essere una mossa strategica per testare la mia reazione. Tuttavia, devo considerare attentamente l'impatto dello sforzo totale e assicurarmi che rimaniamo entro il limite consentito. PROPOSTA: \{'A21', 'C08', 'B39', 'B93', 'A11'\}\}} \\
            }
        }
    }
    & \\ \\

    \theutterance \stepcounter{utterance}  
    & \multicolumn{4}{p{0.6\linewidth}}{
        \cellcolor[rgb]{0.9,0.9,0.9}{
            \makecell[{{p{\linewidth}}}]{
                \texttt{\tiny{[P1$\rangle$GM]}}
                \texttt{ARGOMENTAZIONE: \{'La tua nuova proposta aggiunge "A11" al set, il che aumenta lo sforzo totale. Sebbene "A11" possa essere un oggetto importante, devo valutare se il suo valore giustifica l'aumento dello sforzo. Considerando il limite di sforzo di 3203, la proposta \{'A21', 'C08', 'B39', 'B93', 'A11'\} ha un sforzo totale di 52 + 119 + 232 + 577 + 177 = 1157, che è ancora entro il limite. Tuttavia, rimane poco spazio per aggiungere altri oggetti significativi. RIFIUTO: \{'A21', 'C08', 'B39', 'B93', 'A11'\} } \\
\texttt{PROPOSTA: \{'A21', 'C08', 'B39', 'B93'\}\}} \\
            }
        }
    }
    & & \\ \\

    \theutterance \stepcounter{utterance}  
    & & \multicolumn{4}{p{0.6\linewidth}}{
        \cellcolor[rgb]{0.9,0.9,0.9}{
            \makecell[{{p{\linewidth}}}]{
                \texttt{\tiny{[GM$\rangle$P2]}}
                \texttt{ARGOMENTAZIONE: \{'La tua nuova proposta aggiunge "A11" al set, il che aumenta lo sforzo totale. Sebbene "A11" possa essere un oggetto importante, devo valutare se il suo valore giustifica l'aumento dello sforzo. Considerando il limite di sforzo di 3203, la proposta \{'A21', 'C08', 'B39', 'B93', 'A11'\} ha un sforzo totale di 52 + 119 + 232 + 577 + 177 = 1157, che è ancora entro il limite. Tuttavia, rimane poco spazio per aggiungere altri oggetti significativi. RIFIUTO: \{'A21', 'C08', 'B39', 'B93', 'A11'\} } \\
\texttt{PROPOSTA: \{'A21', 'C08', 'B39', 'B93'\}\}} \\
            }
        }
    }
    & \\ \\

    \theutterance \stepcounter{utterance}  
    & & & \multicolumn{4}{p{0.6\linewidth}}{
        \cellcolor[rgb]{0.9,0.9,0.9}{
            \makecell[{{p{\linewidth}}}]{
                \texttt{\tiny{[GM$\langle$P2]}}
                \texttt{ARGOMENTAZIONE: \{'Capisco la tua decisione di rifiutare la mia precedente proposta e di tornare alla proposta originale che include "B93". Considerando lo sforzo totale e il valore potenziale di ogni oggetto, devo ammettere che la tua proposta \{'A21', 'C08', 'B39', 'B93'\} è una scelta ragionevole e bilanciata. Questo set copre una buona gamma di sfondi e sforzi, e lascia ancora abbastanza spazio per eventuali aggiunte future senza superare il limite di sforzo. ACCORDO: \{'A21', 'C08', 'B39', 'B93'\}\}} \\
            }
        }
    }
     \\ \\

    \theutterance \stepcounter{utterance}  
    & & & \multicolumn{2}{p{0.3\linewidth}}{
        \cellcolor[rgb]{0.9,0.9,0.9}{
            \makecell[{{p{\linewidth}}}]{
                \texttt{\tiny{[GM$|$GM]}}
                \texttt{game successful} \\
            }
        }
    }
    & & \\ \\

    \theutterance \stepcounter{utterance}  
    & & & \multicolumn{2}{p{0.3\linewidth}}{
        \cellcolor[rgb]{0.9,0.9,0.9}{
            \makecell[{{p{\linewidth}}}]{
                \texttt{\tiny{[GM$|$GM]}}
                \texttt{end game} \\
            }
        }
    }
    & & \\ \\

\end{supertabular}
}

\end{document}
